% ********************************************************************
% Two-column project / thesis report
%
% Uses sintr.cls, a repurposed journal class (originally ajstd.cls,
% the ASEAN Journal of Science & Technology for Development) adapted
% for student thesis/project reports. Two-column layout, running
% headers with the report title and author, keyword+abstract block,
% and section numbering suited to a short report rather than a book.
% ********************************************************************
\documentclass[english]{sintr}

% ---- REPORT METADATA (shown in headers/footers) ----
\articletype{Project Report}
\articleshorttile{A Short Title for the Running Header}
\authorheader{J. Researcher}
\submitted{\today}

\title{A Compact Two-Column Template for Thesis or Project Reports}

\author[1]{Jane Researcher}
\affil[1]{Department of Physics, Example Institute}

\keywords{project report, two-column, template}

\abstract{This is a short summary of the report: what was done, the main
method, and the headline result, in three to five sentences. It renders in
the two-column abstract block below the title alongside the keyword list.}

\begin{document}

\maketitle

\section{Introduction}
\label{sec:introduction}

State the problem and motivation here. \texttt{sintr.cls} lays the page out
in two columns with running heads carrying the report title and author, in
the style of a journal article rather than a bound thesis chapter — suited
to a short project report submitted as a single document.

\section{Background}
\label{sec:background}

Summarise related work or the theoretical background the report builds on.
Citations resolve through classic \texttt{bibtex}: for example, this
sentence cites a reference work \cite{knuth1984} to show a resolved
citation rather than an unresolved mark.

\section{Methods}
\label{sec:methods}

\subsection{Setup}

Describe the experimental or computational setup.

\subsection{Analysis}

Describe how results were derived from the setup above.

\section{Results}
\label{sec:results}

Present findings here. Tables and figures use the usual \texttt{tabularx},
\texttt{booktabs} and \texttt{graphicx} machinery, already loaded by the
class.

\section{Conclusion}
\label{sec:conclusion}

Summarise what was learned and, where relevant, what follows from it.

\bibliographystyle{babunsrt}
\bibliography{references}

\end{document}
